\subsection{Spatial Reasoning} \label{sec:spatial-setup}

We ask the LM for the coordinates of objects in a room.
We randomly sample 100 rooms, each with 3 different objects placed in 3 different cardinal directions specified using unit vectors (out of north $(0, -1)$, south $(0, 1)$, east $(1, 0)$, and west $(-1, 0)$ as the default conditions). Though using a downward-facing $y$-axis as the default condition may be counter-intuitive, it is natural when drawing top-to-bottom and is the convention in most image processing libraries such as OpenCV (Python), Pillow (Python), and Processing (Java, JavaScript, Python), as well as graphic design applications such as Adobe Illustrator. We believe this system is the most often encountered during LM pretraining. However, other libraries with an upward-facing $y$-axis also exist, such as matplotlib (Python), ggplot (R), and D3 (JavaScript).

For the counterfactual setting, we alter the direction--unit vector mapping, and ask for the object coordinates in the new system. We consider two direction-swapped worlds (north-south and east-west), three rotated worlds (by 90\textdegree, 180\textdegree, and 270\textdegree), and a randomly permuted world. We evaluate the relative positions of objects and report the instance-level accuracy that requires all 3 objects in a room to be located correctly as the main metric. The random accuracy is around 16.7\%.\footnote{When not considering cases where objects are placed in the same line, there are 24 permutations for placing 3 objects in 4 different directions, of which 4 can be considered correct.} We also report the object-level accuracy in Table~\ref{tab:numbers-spatial}.
As the CCC, we make sure that the LM understands the permuted world by asking it to also specify the coordinates of the unit vectors representing the 4 cardinal directions in the output.
